\section{Data validation}
The dataset used for validating the feasibility of the proposed approach and evaluating the overall functionality of the implementation is the \textit{Electric Motor Vibrations Dataset} provided by the SOON: Social Network of Machines CHIST-ERA Project (\cite{DATA_SET}). The dataset consists of 30 individual column separated values (CSV) files containing vibration measurements collected from two different electrical motors operating under different conditions. The appendix \ref{app:dataset_description} exactly documents each test condition in which the motors where tested in this thesis these different tests will be refereed by their individual ids.
Each dataset contains a timestamp value followed by acceleration measurements along the three spatial axes ($x$, $y$, and $z$). The timestamp column represents the time order of the measurements, while the remaining columns describe the recorded vibration signals of the sensor in the respective axes.
Before the data can be used for anomaly detection, two essential preprocessing and validation steps are required: First, the timestamps must be validated to ensure a consistent sampling interval across the dataset. Second, the three-dimensional acceleration vectors must be transformed into a single representative signal for easier processing.
\subsection{Sampling Interval ($\Delta$t)}
Consistent sampling intervals are important factors in time series data. Large time jumps between consecutive individual measurements can introduce artificial anomalies. Therefore the dataset was validated by calculating the time differences between consecutive timestamps. 
The interval $\Delta t$ was calculated as : \\
\begin{center}
$\Delta t_n = t_n - t_{n-1}$ \\
\end{center}   
Where $\Delta t_n$ is the time  difference between the current  and previous measurement. 
With Python's NumPy min and max functions, the minimum and maximum $\Delta t$ were computed for each dataset. The complete evaluation results are detailed in \ref{app:appendix2}. Inspection identified anomalies characterized by unexpected negative or excessively large values. After excluding these specific files  
(Files: 05, 13, 14, and 16) and running again the algorithm, the resulting statistical values improved significantly, as show in \ref{table:statistics}.\\
\begin{table}[!h]
 \centering
 \begin{tabular}{|c|c|}
 \hline
 Total number of intervals & 2,975,473 \\ 
 \hline
 Mean $\Delta t$ & 1,906.94 $\mu$s \\ 
 \hline
 Standard deviation & 38.03 $\mu$s \\ 
 \hline
 Median & 1,912.0 $\mu$s \\
 \hline
 Minimum & 1,660.0 $\mu$s \\
 \hline
 Maximum & 2,032.0 $\mu$s \\
 \hline
 \end{tabular}
 \caption{Statistics of sampling intervals after filtering rounded to the second decimal}
 \label{table:statistics}
\end{table}
\\
By analyzing the results, a clear improvement can be observed, the difference between the minimum, maximum, and the global mean are very small. Additionally, the standard deviation has improved a lot, which indicates highly consistent sampling. These results imply that the dataset was successful cleaned and can be used for further benchmarking and testing.
\subsection{Combining vectors}
A method for combining three orthogonal acceleration vectors into a single composite measure is the vector magnitude approach, as it preserves the total acceleration energy of the signal. By only taking the sum or average of the three axes bias might be introduced or important information might be lost. To obtain a single measure of total acceleration, the vector magnitude is calculated for each dataset using the theorem:
\begin{equation}
    VM = \sqrt{(x^2 + y^2 + z^2)}
\end{equation}
This method is supported by \cite{formulaCombiningVectors}, who state that "the VM equation is the square root of the sum of the squares of three planes of motion."